\scanpage{62}
\mspage{\unsure{48}}
% archive sheet number: 61

où l'objet $Z_{\unsure{i_0|0}}$, et les trois flèches $\alpha_Z,i_0,p_0$ sont à déterminer. Si on a $i'\in C_{I'^{\unsure{\circ|0}}},p\in W_{I'^\circ}$, alors on conclut par th.~2.1 $1^\circ$ (\unsure{cor.} 2.2) que \inserted{$i'_{I'}\in C$,} et si on veut que $i\in C_{I^\circ},p\in W_{I^\circ}$, il reste seulement à exiger \deleted{d'une part que}

\[
 p_0\in W,\ (i'_{I',I}:\underbrace{Z_{I'}\amalg_{X_{I'}}X_0}_{Z'_0}\longrightarrow Z_0)\in C
\]

d'après le th.~2.1 $2^\circ$. On va alors construire le diagramme $(*)$, à partir du diagramme similaire ci-contre (avec $Z_0$ de $(*)$ remplacé par $Z'_0=Z_{I',I}$)

\diagram{\begin{tikzcd}[ampersand replacement=\&,column sep=small,row sep=small]
 \& X_{I'}\arrow[rr,"\alpha_X"]\arrow[dl,"{i'_{I'}\in C}"']\arrow[dd,"{f'_{I'}}"] \& \& X_0\arrow[dl,"{i'_0}"']\arrow[dd,"{f_0}"]\arrow[dr,dashed,"{i_0=j_0i'_0}"] \&\\
 Z_{I'}\arrow[rr,"{\alpha'_Z}"]\arrow[dr,"{p'_{I'}}"']\arrow[urrr,dotted,leftrightarrow,shorten <=1.6em,shorten >=2.7em,"{\text{cocart}}" description,sloped] \& \& Z'_0=Z_{I',I}\arrow[rr,dash,dotted,very thick]\arrow[rr,thick,"{j_0\in C}"]\arrow[dr,"{p'_0}"'] \& \& Z_0\arrow[dl,dashed,"{p_0\in W}"]\\
 \& Y_{I'}\arrow[rr,"\alpha_Y"] \& \& Y_0 \&
\end{tikzcd}}

en factorisant la flèche canonique $p'_0:Z'_0\longrightarrow Y_0$ en \deleted{$\unsure{p}$}
\[
 p'_0=\deleted{\unsure{q}}p_0j_0,\quad\text{avec}\quad \deleted{\unsure{q}}p_0\in W,\ j_0\in C,
\]

et en posant
\[
 i_0=j_0i'_0,
\]
ce qui achève la construction.

On a ainsi prouvé le
